\documentclass[10pt,a4paper]{amsart}
\usepackage[margin=19mm]{geometry}
\usepackage{amsmath,amssymb,mathtools}
\usepackage[scr=rsfs,cal=euler]{mathalfa}
\usepackage{xfrac}
\usepackage[unicode,hidelinks,bookmarks=false]{hyperref}
\newcommand{\ie}{i.e.\ }
\newcommand{\cf}{cf.\ }
\newcommand{\resp}{respectively\ }
\newcommand{\Subsection}{\subsection}
\providecommand{\nobreakdash}{\nobreak\hbox{-}}
\setlength{\emergencystretch}{2em}
\multlinegap0pt
\usepackage[shortlabels]{enumitem}
\usepackage{mathtools}
\usepackage{tensor}
\usepackage{float}
\usepackage{amssymb}
\usepackage{tikz}
\usepackage{csquotes}
\newtheorem{theorem}{Theorem}
\title{$\lambda$-Chromatic Polynomials and Polytope Geometry}
\author{Annayat Ali, Rameez Raja}
\date{Source: arXiv:2605.31047v1 (CC BY 4.0)}
\begin{document}
\maketitle

\noindent\emph{Source and licence.} $\lambda$-Chromatic Polynomials and Polytope Geometry. Version arXiv:2605.31047v1 (CC BY 4.0), distributed by arXiv under the Creative Commons Attribution 4.0 International licence, http://creativecommons.org/licenses/by/4.0/, as recorded in the arXiv licence metadata for this exact version and shown on the abstract page https://arxiv.org/abs/2605.31047v1. Full licence notice: \texttt{licenses/arxiv-2605.31047-CC-BY-4.0.txt}.\par\smallskip

\noindent\textit{Editorial scope.} Counted unit: the article's theorem LKN (source0.tex lines 520--526). The article's complete original proof of it is delivered with it (source0.tex lines 527--577, one \texttt{proof} environment of 51 lines). No mathematics is written, repaired or re-derived: the statement and the proof are the article's own lines, in the article's own notation, and no retained line is substituted or re-flowed.  Retained uncounted as the source-local context the proof needs: lines 388--391.  Nothing was omitted from any retained span, so the package carries no omission record.  No retained line contains a citation, so the wrapper carries no bibliography and no bibliography entry.  No retained line contains a figure environment, an image-inclusion command or drawing code, so no image is delivered and the package declares no asset dependency.  Editorial formatting only, in addition: an AMS \texttt{amsart} preamble that restates the article's own declarations and macros used by the retained lines, namely the environments \texttt{theorem} on separate counters, together with the standard packages those lines need.  The retained environments print in this document as Theorem 1 in order of appearance, whereas the article prints the same items with the numbering of its own full document.  The complete native source of the article as distributed in the arXiv e-print for this exact version is bundled unchanged as \texttt{sources/201705/source0.tex}.
\begin{theorem}
\label{t1}
For the complete graph \( K_n \), there is an one-to-one correspondence between valid \( L(2,1) \)-colourings of \( K_n \) with maximum colour at most \( m \) and North-East lattice paths in \( \mathcal{C}(K_n, f, m) \) that begin in the bottom row and terminate in the penultimate row.
\end{theorem}

\begin{theorem}
\label{LKN}
The \( \lambda \)-chromatic polynomial of the complete graph \( K_n \) is given by
\[
\Lambda(K_n, x) = \ (x - n + 2)(x - n + 1)(x-n) \cdots (x - 2n + 3).
\]
\end{theorem}

\begin{proof}
The $\lambda$-colouring of the complete graph \( K_n \) is given by the function
\[
f(v_i) = 2i, \quad \text{where } 0 \leq i \leq n-1,
\]
where the vertex set is \( V(K_n) = \{v_0, v_1, \dots, v_{n-1}\} \). This is the unique $\lambda$-colouring of \( K_n \).

The problem of counting distinct $L(2,1)$-colourings using the colour set \( \{0,1, \dots, m\} \) reduces to counting the number of north-east lattice paths in the colouring grid \( \mathcal{C}(K_n, f, m) \), as suggested by Theorem~\ref{t1}.

More precisely, we count the number of north-east lattice paths from a cell \( (a,0) \) to \( (b,n-2) \), where \( 0 \leq a \leq x - 2n + 2 \) and \( a \leq b \leq x - 2n + 2 \).

Substitute \( N = n - 2 \) and \( M = x - 2n + 2 \). It is well known that the number of north-east lattice paths from \( (0,0) \) to \( (a,b) \) is given by the binomial coefficient:
\[
\binom{a+b}{a}.
\]
First, we count the number of ways to travel from a fixed starting cell \( (i,0) \) to a fixed ending cell \( (j,N) \), where \( 0 \leq i \leq j \leq M \). This count is given by:
\[
\binom{N+(j-i)}{N}.
\]
Summing over all possible paths starting from any cell in the bottom row \( (i,0) \) and ending at any cell in the top row \( (j,N) \), we obtain:
\[
\sum\limits_{i=0}^{M} \sum\limits_{j=i}^{M} \binom{N+(j-i)}{N}.
\]
Now, setting \( k = j - i \), we rewrite the sum as:
\[
\sum\limits_{i=0}^{M} \sum\limits_{k=0}^{M-i} \binom{N+k}{N}.
\]
Applying the hockey-stick identity,
\[
\sum\limits_{k=0}^{m} \binom{N+k}{N} = \binom{N+m+1}{N+1},
\]
we simplify the inner sum to:
\[
\binom{N+(M-i)+1}{N+1}.
\]
Now, summing over all \( i \) from \( 0 \) to \( M \), we again apply the hockey-stick identity:
\[
\sum\limits_{i=0}^{M} \binom{N+(M-i)+1}{N+1} = \binom{N+M+2}{N+2}.
\]
Finally, substituting back \( M = x- 2n + 2 \) and \( N = n - 2 \), we obtain:
\[
\binom{x-n+2}{n}.
\]
Given an $L(2,1)$-colouring of \( K_n \), we can permute the colours, and the colouring remains valid. Since there are \( n! \) ways to permute the colours, we multiply by \( n! \) to account for these permutations. Therefore,

\[
\Lambda(K_n, x) = (x-n+2)(x-n+1)\cdots (x-2n+3).
\]

 This completes the proof.
\end{proof}

\end{document}
